% BHU PYQ -- Manually transcribed & error-corrected
% Paper: BCH-222 : Business Economics-II
% Exam: B.Com. (Hons.) Semester IV Examination, 2023-24

\documentclass[11pt,a4paper]{article}
\usepackage[a4paper,top=2.5cm,bottom=2.5cm,left=2.8cm,right=2.8cm,headheight=14pt]{geometry}
\usepackage{amsmath,amssymb}
\usepackage[T1]{fontenc}
\usepackage[utf8]{inputenc}
\usepackage{lmodern,microtype,fancyhdr,enumitem,hyperref,lastpage,parskip}

\hypersetup{
    colorlinks=true,
    urlcolor=black,
    linkcolor=black,
    pdftitle={BCH-222: Business Economics-II -- BHU B.Com. (Hons.) Sem IV 2023-24}
}

\pagestyle{fancy}
\fancyhf{}
\renewcommand{\headrulewidth}{0.4pt}
\renewcommand{\footrulewidth}{0.4pt}
\fancyhead[L]{\small   \textit{Banaras Hindu University}}
\fancyhead[R]{\small   \textit{BCH-222: Business Economics-II}}
\fancyfoot[C]{\small Page  \thepage\ of \pageref{LastPage}}


\newlist{parts}{enumerate}{2}
\setlist[parts,1]{label=(\roman*),leftmargin=*,itemsep=5pt,topsep=3pt}
\setlist[parts,2]{label=(\alph*),leftmargin=*,itemsep=3pt,topsep=2pt}

\newcommand{\pts}[1]{\hfill   \textit{[#1]}}


\begin{document}


\begin{center}
  {\large\bfseries BANARAS HINDU UNIVERSITY}\\[2pt]
  {\normalsize B.Com. (Hons.) --- Semester IV Examination, 2023-24}\\[6pt]
  \rule{0.8\linewidth}{0.5pt}\\[6pt]
  {\large\bfseries Paper No. BCH-222: Business Economics-II}\\[4pt]
  {\normalsize Time: 3 Hours\hfill Maximum Marks: 70}
\end{center}

\rule{\linewidth}{0.4pt}

\smallskip



\noindent   \textit{Instructions: Answer all questions. All questions carry equal marks (14 marks each).}

\bigskip




\noindent   \textbf{Question 1.}
Explain the different types of Market structure and describe the conditions for the firm's equilibrium under perfect competition.\pts{14}

\centerline{   \textbf{OR}}

What do you understand by the term `Monopoly'? Explain the price-output relationships in discriminating monopoly market structure.\pts{14}

\medskip\rule{\linewidth}{0.2pt}




\noindent   \textbf{Question 2.}
Explain the features of oligopoly. Also discuss its types.\pts{14}

\centerline{   \textbf{OR}}

How does a firm determine its equilibrium both in short period and long period under monopolistic competition?\pts{14}

\medskip\rule{\linewidth}{0.2pt}




\noindent   \textbf{Question 3.}
Explain the Ricardian Theory of Rent with proper illustration and examples.\pts{14}

\centerline{   \textbf{OR}}

Comment on the following statement, ``Profits result from societal changes that are dynamic in nature.''\pts{14}

\medskip\rule{\linewidth}{0.2pt}




\noindent   \textbf{Question 4.}
Discuss the Marginal Productivity Theory of Distribution along with its assumptions.\pts{14}

\centerline{   \textbf{OR}}

Critically explain the marginal productivity theory of wages. Contrast this theory with the modern theory of wages.\pts{14}

\medskip\rule{\linewidth}{0.2pt}




\noindent   \textbf{Question 5.}
Write short notes on any two of the following:

\begin{parts}
  \item Risk bearing theory of Profit
  \item Keynesian theory of interest
  \item Quasi-Rent
\end{parts}\pts{14}

\centerline{   \textbf{OR}}

Write short notes on the following:

\begin{parts}
  \item Product Differentiation
  \item Gross Profit and Normal Profit
\end{parts}\pts{14}

\vfill
\rule{\linewidth}{0.4pt}

\begin{center}\small
\href{https://bhu-exam.vercel.app/}{Click here to visit our website}\\[4pt]
\href{https://me-aryan.vercel.app/}{Click here to visit our contributor}\\[4pt]
\href{https://quashan-me.vercel.app/}{Click here to visit our contributor}
\end{center}

\end{document}
